% !PRINTED 552
\weberpage{552}
%<<CONT
be, where $a$ is a constant (independent of $t$), and $\theta$ a function of $t$ which remains enclosed within finite limits.

If we assume first that the $\mu_n$ are all different from one another, then $Z(\mu_n) = n$, and from (2) it follows that:
\weq{3}{\frac{n}{\mu_n} = a + \frac{\theta}{\sqrt{\mu_n}};}
from this it follows that $n:\mu_n$ remains finite, and consequently, with an altered meaning of $\theta$:
\[
\frac{n}{\mu_n} = a + \frac{\theta}{\sqrt{n}},
\]
or, by the binomial theorem, for any positive $s$:
\weq{4}{\frac{n^s}{\mu_n^s} = a^s\left(1 + \frac{\theta}{\sqrt{n}}\right),}
in which all the quantities denoted by $\theta$ have finite values.

But these formulae still hold even when equal ones occur among the $\mu_n$. For let, say (as in Bd.\,II, § 196, 3.)
\[
\mu_{m+1}, \mu_{m+2}, \ldots, \mu_{m+l} = \mu_n
\]
be equal to one another and
\[
m < n \leqq m+l,
\]
then
\[
Z(\mu_n-0) = m, \qquad Z(\mu_n) = m+l,
\]
\[
\frac{Z(\mu_n-0)}{\mu_n} < \frac{n}{\mu_n} \leqq \frac{Z(\mu_n)}{\mu_n},
\]
\[
a + \frac{\theta'}{\sqrt{\mu_n}} < \frac{n}{\mu_n} < a + \frac{\theta}{\sqrt{\mu_n}},
\]
from which (3) follows and (4) can be derived as above.

The theorem we wish to prove runs as follows:

\gesperrt{Let $\varepsilon_1, \varepsilon_2, \varepsilon_3, \ldots$ be an infinite sequence of positive or negative, but finite quantities and}
\weq{5}{\varepsilon_1+\varepsilon_2+\cdots+\varepsilon_n = \gamma_n\sqrt{n},}
\gesperrt{so constituted that $\gamma_n$ remains, in absolute value, below a finite constant.}

\gesperrt{Then}
\weq{6}{\sigma = \frac{\varepsilon_1}{\mu_1^s} + \frac{\varepsilon_2}{\mu_2^s} + \frac{\varepsilon_3}{\mu_3^s} + \cdots}
\gesperrt{is convergent and a continuous function of $s$, so long as}
\[
s > \tfrac{1}{2}
\]
\gesperrt{holds.}
% [sic] Printed binomial step(4) requires a>0, whereas the stated hypothesis does not exclude a=0. Formula retained.
